% Options for packages loaded elsewhere
\PassOptionsToPackage{unicode}{hyperref}
\PassOptionsToPackage{hyphens}{url}
\documentclass[
]{article}
\usepackage{xcolor}
\usepackage[margin=1in]{geometry}
\usepackage{amsmath,amssymb}
\setcounter{secnumdepth}{-\maxdimen} % remove section numbering
\usepackage{iftex}
\ifPDFTeX
  \usepackage[T1]{fontenc}
  \usepackage[utf8]{inputenc}
  \usepackage{textcomp} % provide euro and other symbols
\else % if luatex or xetex
  \usepackage{unicode-math} % this also loads fontspec
  \defaultfontfeatures{Scale=MatchLowercase}
  \defaultfontfeatures[\rmfamily]{Ligatures=TeX,Scale=1}
\fi
\usepackage{lmodern}
\ifPDFTeX\else
  % xetex/luatex font selection
\fi
% Use upquote if available, for straight quotes in verbatim environments
\IfFileExists{upquote.sty}{\usepackage{upquote}}{}
\IfFileExists{microtype.sty}{% use microtype if available
  \usepackage[]{microtype}
  \UseMicrotypeSet[protrusion]{basicmath} % disable protrusion for tt fonts
}{}
\makeatletter
\@ifundefined{KOMAClassName}{% if non-KOMA class
  \IfFileExists{parskip.sty}{%
    \usepackage{parskip}
  }{% else
    \setlength{\parindent}{0pt}
    \setlength{\parskip}{6pt plus 2pt minus 1pt}}
}{% if KOMA class
  \KOMAoptions{parskip=half}}
\makeatother
\usepackage{color}
\usepackage{fancyvrb}
\newcommand{\VerbBar}{|}
\newcommand{\VERB}{\Verb[commandchars=\\\{\}]}
\DefineVerbatimEnvironment{Highlighting}{Verbatim}{commandchars=\\\{\}}
% Add ',fontsize=\small' for more characters per line
\usepackage{framed}
\definecolor{shadecolor}{RGB}{248,248,248}
\newenvironment{Shaded}{\begin{snugshade}}{\end{snugshade}}
\newcommand{\AlertTok}[1]{\textcolor[rgb]{0.94,0.16,0.16}{#1}}
\newcommand{\AnnotationTok}[1]{\textcolor[rgb]{0.56,0.35,0.01}{\textbf{\textit{#1}}}}
\newcommand{\AttributeTok}[1]{\textcolor[rgb]{0.13,0.29,0.53}{#1}}
\newcommand{\BaseNTok}[1]{\textcolor[rgb]{0.00,0.00,0.81}{#1}}
\newcommand{\BuiltInTok}[1]{#1}
\newcommand{\CharTok}[1]{\textcolor[rgb]{0.31,0.60,0.02}{#1}}
\newcommand{\CommentTok}[1]{\textcolor[rgb]{0.56,0.35,0.01}{\textit{#1}}}
\newcommand{\CommentVarTok}[1]{\textcolor[rgb]{0.56,0.35,0.01}{\textbf{\textit{#1}}}}
\newcommand{\ConstantTok}[1]{\textcolor[rgb]{0.56,0.35,0.01}{#1}}
\newcommand{\ControlFlowTok}[1]{\textcolor[rgb]{0.13,0.29,0.53}{\textbf{#1}}}
\newcommand{\DataTypeTok}[1]{\textcolor[rgb]{0.13,0.29,0.53}{#1}}
\newcommand{\DecValTok}[1]{\textcolor[rgb]{0.00,0.00,0.81}{#1}}
\newcommand{\DocumentationTok}[1]{\textcolor[rgb]{0.56,0.35,0.01}{\textbf{\textit{#1}}}}
\newcommand{\ErrorTok}[1]{\textcolor[rgb]{0.64,0.00,0.00}{\textbf{#1}}}
\newcommand{\ExtensionTok}[1]{#1}
\newcommand{\FloatTok}[1]{\textcolor[rgb]{0.00,0.00,0.81}{#1}}
\newcommand{\FunctionTok}[1]{\textcolor[rgb]{0.13,0.29,0.53}{\textbf{#1}}}
\newcommand{\ImportTok}[1]{#1}
\newcommand{\InformationTok}[1]{\textcolor[rgb]{0.56,0.35,0.01}{\textbf{\textit{#1}}}}
\newcommand{\KeywordTok}[1]{\textcolor[rgb]{0.13,0.29,0.53}{\textbf{#1}}}
\newcommand{\NormalTok}[1]{#1}
\newcommand{\OperatorTok}[1]{\textcolor[rgb]{0.81,0.36,0.00}{\textbf{#1}}}
\newcommand{\OtherTok}[1]{\textcolor[rgb]{0.56,0.35,0.01}{#1}}
\newcommand{\PreprocessorTok}[1]{\textcolor[rgb]{0.56,0.35,0.01}{\textit{#1}}}
\newcommand{\RegionMarkerTok}[1]{#1}
\newcommand{\SpecialCharTok}[1]{\textcolor[rgb]{0.81,0.36,0.00}{\textbf{#1}}}
\newcommand{\SpecialStringTok}[1]{\textcolor[rgb]{0.31,0.60,0.02}{#1}}
\newcommand{\StringTok}[1]{\textcolor[rgb]{0.31,0.60,0.02}{#1}}
\newcommand{\VariableTok}[1]{\textcolor[rgb]{0.00,0.00,0.00}{#1}}
\newcommand{\VerbatimStringTok}[1]{\textcolor[rgb]{0.31,0.60,0.02}{#1}}
\newcommand{\WarningTok}[1]{\textcolor[rgb]{0.56,0.35,0.01}{\textbf{\textit{#1}}}}
\usepackage{graphicx}
\makeatletter
\newsavebox\pandoc@box
\newcommand*\pandocbounded[1]{% scales image to fit in text height/width
  \sbox\pandoc@box{#1}%
  \Gscale@div\@tempa{\textheight}{\dimexpr\ht\pandoc@box+\dp\pandoc@box\relax}%
  \Gscale@div\@tempb{\linewidth}{\wd\pandoc@box}%
  \ifdim\@tempb\p@<\@tempa\p@\let\@tempa\@tempb\fi% select the smaller of both
  \ifdim\@tempa\p@<\p@\scalebox{\@tempa}{\usebox\pandoc@box}%
  \else\usebox{\pandoc@box}%
  \fi%
}
% Set default figure placement to htbp
\def\fps@figure{htbp}
\makeatother
\setlength{\emergencystretch}{3em} % prevent overfull lines
\providecommand{\tightlist}{%
  \setlength{\itemsep}{0pt}\setlength{\parskip}{0pt}}
\usepackage{bookmark}
\IfFileExists{xurl.sty}{\usepackage{xurl}}{} % add URL line breaks if available
\urlstyle{same}
\hypersetup{
  pdftitle={Proyecto de Probabilidad y Estadística},
  pdfauthor={Autor},
  hidelinks,
  pdfcreator={LaTeX via pandoc}}

\title{Proyecto de Probabilidad y Estadística}
\author{Autor}
\date{2026-04-14}

\usepackage{graphicx}
\graphicspath{{Tallerprobabilidad/figure-latex/}}

\begin{document}
\maketitle

{
\setcounter{tocdepth}{3}
\tableofcontents
}
\section{Introducción}\label{introducciuxf3n}

Este proyecto aborda dos problemas fundamentales de probabilidad:

\begin{enumerate}
\def\labelenumi{\arabic{enumi}.}
\tightlist
\item
  Selección de estudiantes (combinatoria)
\item
  Lanzamiento de dos dados (variable aleatoria discreta)
\end{enumerate}

Se analizan tanto desde un enfoque teórico como mediante simulación.

\section{PARTE 1: COMBINATORIA Y
PROBABILIDAD}\label{parte-1-combinatoria-y-probabilidad}

\subsection{Estructura de datos}\label{estructura-de-datos}

\begin{Shaded}
\begin{Highlighting}[]
\NormalTok{mujeres }\OtherTok{\textless{}{-}} \FunctionTok{c}\NormalTok{(}\StringTok{"MA"}\NormalTok{, }\StringTok{"MB"}\NormalTok{, }\StringTok{"MC"}\NormalTok{, }\StringTok{"MD"}\NormalTok{)}
\NormalTok{hombres }\OtherTok{\textless{}{-}} \FunctionTok{c}\NormalTok{(}\StringTok{"HA"}\NormalTok{, }\StringTok{"HB"}\NormalTok{, }\StringTok{"HC"}\NormalTok{, }\StringTok{"HD"}\NormalTok{, }\StringTok{"HE"}\NormalTok{, }\StringTok{"HF"}\NormalTok{)}
\NormalTok{estudiantes }\OtherTok{\textless{}{-}} \FunctionTok{c}\NormalTok{(mujeres, hombres)}

\FunctionTok{length}\NormalTok{(estudiantes)}
\end{Highlighting}
\end{Shaded}

\begin{verbatim}
## [1] 10
\end{verbatim}

\subsection{Espacio muestral}\label{espacio-muestral}

\begin{Shaded}
\begin{Highlighting}[]
\NormalTok{espacio\_muestral }\OtherTok{\textless{}{-}} \FunctionTok{combn}\NormalTok{(estudiantes, }\DecValTok{3}\NormalTok{, }\AttributeTok{simplify =} \ConstantTok{FALSE}\NormalTok{)}
\NormalTok{espacio\_muestral\_str }\OtherTok{\textless{}{-}} \FunctionTok{sapply}\NormalTok{(espacio\_muestral, }\ControlFlowTok{function}\NormalTok{(x) }\FunctionTok{paste}\NormalTok{(}\FunctionTok{sort}\NormalTok{(x), }\AttributeTok{collapse =} \StringTok{"{-}"}\NormalTok{))}

\FunctionTok{length}\NormalTok{(espacio\_muestral)}
\end{Highlighting}
\end{Shaded}

\begin{verbatim}
## [1] 120
\end{verbatim}

\begin{Shaded}
\begin{Highlighting}[]
\FunctionTok{head}\NormalTok{(espacio\_muestral\_str)}
\end{Highlighting}
\end{Shaded}

\begin{verbatim}
## [1] "MA-MB-MC" "MA-MB-MD" "HA-MA-MB" "HB-MA-MB" "HC-MA-MB" "HD-MA-MB"
\end{verbatim}

\subsection{Funciones auxiliares}\label{funciones-auxiliares}

\begin{Shaded}
\begin{Highlighting}[]
\NormalTok{contar\_mujeres }\OtherTok{\textless{}{-}} \ControlFlowTok{function}\NormalTok{(muestra) \{}
  \FunctionTok{sum}\NormalTok{(muestra }\SpecialCharTok{\%in\%}\NormalTok{ mujeres)}
\NormalTok{\}}

\NormalTok{contiene\_A\_y\_D }\OtherTok{\textless{}{-}} \ControlFlowTok{function}\NormalTok{(muestra) \{}
  \StringTok{"MA"} \SpecialCharTok{\%in\%}\NormalTok{ muestra }\SpecialCharTok{\&\&} \StringTok{"MD"} \SpecialCharTok{\%in\%}\NormalTok{ muestra}
\NormalTok{\}}
\end{Highlighting}
\end{Shaded}

\subsection{Clasificación}\label{clasificaciuxf3n}

\begin{Shaded}
\begin{Highlighting}[]
\NormalTok{clasificacion }\OtherTok{\textless{}{-}} \FunctionTok{data.frame}\NormalTok{(}
  \AttributeTok{muestra =}\NormalTok{ espacio\_muestral\_str,}
  \AttributeTok{num\_mujeres =} \FunctionTok{sapply}\NormalTok{(espacio\_muestral, contar\_mujeres),}
  \AttributeTok{tiene\_A\_y\_D =} \FunctionTok{sapply}\NormalTok{(espacio\_muestral, contiene\_A\_y\_D)}
\NormalTok{)}

\FunctionTok{head}\NormalTok{(clasificacion)}
\end{Highlighting}
\end{Shaded}

\begin{verbatim}
##    muestra num_mujeres tiene_A_y_D
## 1 MA-MB-MC           3       FALSE
## 2 MA-MB-MD           3        TRUE
## 3 HA-MA-MB           2       FALSE
## 4 HB-MA-MB           2       FALSE
## 5 HC-MA-MB           2       FALSE
## 6 HD-MA-MB           2       FALSE
\end{verbatim}

\subsection{Probabilidades}\label{probabilidades}

\begin{Shaded}
\begin{Highlighting}[]
\NormalTok{prob\_dos\_mujeres }\OtherTok{\textless{}{-}} \FunctionTok{mean}\NormalTok{(clasificacion}\SpecialCharTok{$}\NormalTok{num\_mujeres }\SpecialCharTok{==} \DecValTok{2}\NormalTok{)}
\NormalTok{prob\_al\_menos\_una }\OtherTok{\textless{}{-}} \FunctionTok{mean}\NormalTok{(clasificacion}\SpecialCharTok{$}\NormalTok{num\_mujeres }\SpecialCharTok{\textgreater{}=} \DecValTok{1}\NormalTok{)}
\NormalTok{prob\_cero\_mujeres }\OtherTok{\textless{}{-}} \FunctionTok{mean}\NormalTok{(clasificacion}\SpecialCharTok{$}\NormalTok{num\_mujeres }\SpecialCharTok{==} \DecValTok{0}\NormalTok{)}
\NormalTok{prob\_sin\_A\_y\_D }\OtherTok{\textless{}{-}} \FunctionTok{mean}\NormalTok{(}\SpecialCharTok{!}\NormalTok{clasificacion}\SpecialCharTok{$}\NormalTok{tiene\_A\_y\_D)}

\FunctionTok{c}\NormalTok{(prob\_dos\_mujeres,}
\NormalTok{  prob\_al\_menos\_una,}
\NormalTok{  prob\_cero\_mujeres,}
\NormalTok{  prob\_sin\_A\_y\_D)}
\end{Highlighting}
\end{Shaded}

\begin{verbatim}
## [1] 0.3000000 0.8333333 0.1666667 0.9333333
\end{verbatim}

\subsection{Resumen}\label{resumen}

\begin{Shaded}
\begin{Highlighting}[]
\NormalTok{resumen1 }\OtherTok{\textless{}{-}} \FunctionTok{data.frame}\NormalTok{(}
  \AttributeTok{Condicion =} \FunctionTok{c}\NormalTok{(}\StringTok{"Exactamente 2 mujeres"}\NormalTok{, }
                \StringTok{"Al menos 1 mujer"}\NormalTok{, }
                \StringTok{"Sin mujeres"}\NormalTok{, }
                \StringTok{"MA y MD NO juntas"}\NormalTok{),}
  \AttributeTok{Probabilidad =} \FunctionTok{c}\NormalTok{(prob\_dos\_mujeres, prob\_al\_menos\_una, prob\_cero\_mujeres, prob\_sin\_A\_y\_D)}
\NormalTok{)}

\NormalTok{resumen1}
\end{Highlighting}
\end{Shaded}

\begin{verbatim}
##               Condicion Probabilidad
## 1 Exactamente 2 mujeres    0.3000000
## 2      Al menos 1 mujer    0.8333333
## 3           Sin mujeres    0.1666667
## 4     MA y MD NO juntas    0.9333333
\end{verbatim}

\subsection{Simulación}\label{simulaciuxf3n}

\begin{Shaded}
\begin{Highlighting}[]
\FunctionTok{set.seed}\NormalTok{(}\DecValTok{123}\NormalTok{)}
\NormalTok{n\_sim }\OtherTok{\textless{}{-}} \DecValTok{100000}

\NormalTok{resultados\_sim }\OtherTok{\textless{}{-}} \FunctionTok{replicate}\NormalTok{(n\_sim, }\FunctionTok{sample}\NormalTok{(estudiantes, }\DecValTok{3}\NormalTok{))}

\FunctionTok{mean}\NormalTok{(}\FunctionTok{apply}\NormalTok{(resultados\_sim, }\DecValTok{2}\NormalTok{, contar\_mujeres) }\SpecialCharTok{==} \DecValTok{2}\NormalTok{)}
\end{Highlighting}
\end{Shaded}

\begin{verbatim}
## [1] 0.30141
\end{verbatim}

\subsection{Visualización}\label{visualizaciuxf3n}

\begin{Shaded}
\begin{Highlighting}[]
\NormalTok{tabla }\OtherTok{\textless{}{-}} \FunctionTok{table}\NormalTok{(clasificacion}\SpecialCharTok{$}\NormalTok{num\_mujeres)}

\FunctionTok{barplot}\NormalTok{(tabla,}
        \AttributeTok{main =} \StringTok{"Distribución del número de mujeres"}\NormalTok{,}
        \AttributeTok{xlab =} \StringTok{"Número de mujeres"}\NormalTok{,}
        \AttributeTok{ylab =} \StringTok{"Frecuencia"}\NormalTok{)}
\end{Highlighting}
\end{Shaded}

\pandocbounded{\includegraphics[keepaspectratio]{Tallerprobabilidad/figure-latex/unnamed-chunk-8-1.pdf}}

\section{PARTE 2: LANZAMIENTO DE DOS
DADOS}\label{parte-2-lanzamiento-de-dos-dados}

\subsection{Espacio muestral}\label{espacio-muestral-1}

\begin{Shaded}
\begin{Highlighting}[]
\NormalTok{dado1 }\OtherTok{\textless{}{-}} \DecValTok{1}\SpecialCharTok{:}\DecValTok{6}
\NormalTok{dado2 }\OtherTok{\textless{}{-}} \DecValTok{1}\SpecialCharTok{:}\DecValTok{6}

\NormalTok{espacio\_muestral\_dados }\OtherTok{\textless{}{-}} \FunctionTok{expand.grid}\NormalTok{(}\AttributeTok{Dado1 =}\NormalTok{ dado1, }\AttributeTok{Dado2 =}\NormalTok{ dado2)}
\NormalTok{espacio\_muestral\_dados}\SpecialCharTok{$}\NormalTok{Suma }\OtherTok{\textless{}{-}}\NormalTok{ espacio\_muestral\_dados}\SpecialCharTok{$}\NormalTok{Dado1 }\SpecialCharTok{+}\NormalTok{ espacio\_muestral\_dados}\SpecialCharTok{$}\NormalTok{Dado2}

\FunctionTok{head}\NormalTok{(espacio\_muestral\_dados)}
\end{Highlighting}
\end{Shaded}

\begin{verbatim}
##   Dado1 Dado2 Suma
## 1     1     1    2
## 2     2     1    3
## 3     3     1    4
## 4     4     1    5
## 5     5     1    6
## 6     6     1    7
\end{verbatim}

\subsection{Distribución de
probabilidad}\label{distribuciuxf3n-de-probabilidad}

\begin{Shaded}
\begin{Highlighting}[]
\NormalTok{tabla\_frecuencias }\OtherTok{\textless{}{-}} \FunctionTok{table}\NormalTok{(espacio\_muestral\_dados}\SpecialCharTok{$}\NormalTok{Suma)}
\NormalTok{valores\_W }\OtherTok{\textless{}{-}} \FunctionTok{as.numeric}\NormalTok{(}\FunctionTok{names}\NormalTok{(tabla\_frecuencias))}
\NormalTok{frecuencias }\OtherTok{\textless{}{-}} \FunctionTok{as.numeric}\NormalTok{(tabla\_frecuencias)}

\NormalTok{probabilidades }\OtherTok{\textless{}{-}}\NormalTok{ frecuencias }\SpecialCharTok{/} \FunctionTok{sum}\NormalTok{(frecuencias)}
\NormalTok{prob\_acumulada }\OtherTok{\textless{}{-}} \FunctionTok{cumsum}\NormalTok{(probabilidades)}

\NormalTok{dist\_W }\OtherTok{\textless{}{-}} \FunctionTok{data.frame}\NormalTok{(}
  \AttributeTok{W =}\NormalTok{ valores\_W,}
  \AttributeTok{Probabilidad =}\NormalTok{ probabilidades,}
  \AttributeTok{Acumulada =}\NormalTok{ prob\_acumulada}
\NormalTok{)}

\NormalTok{dist\_W}
\end{Highlighting}
\end{Shaded}

\begin{verbatim}
##     W Probabilidad  Acumulada
## 1   2   0.02777778 0.02777778
## 2   3   0.05555556 0.08333333
## 3   4   0.08333333 0.16666667
## 4   5   0.11111111 0.27777778
## 5   6   0.13888889 0.41666667
## 6   7   0.16666667 0.58333333
## 7   8   0.13888889 0.72222222
## 8   9   0.11111111 0.83333333
## 9  10   0.08333333 0.91666667
## 10 11   0.05555556 0.97222222
## 11 12   0.02777778 1.00000000
\end{verbatim}

\subsection{Visualizaciones}\label{visualizaciones}

\begin{Shaded}
\begin{Highlighting}[]
\FunctionTok{barplot}\NormalTok{(probabilidades,}
        \AttributeTok{names.arg =}\NormalTok{ valores\_W,}
        \AttributeTok{main =} \StringTok{"Función de Probabilidad P(W=w)"}\NormalTok{,}
        \AttributeTok{xlab =} \StringTok{"Suma"}\NormalTok{,}
        \AttributeTok{ylab =} \StringTok{"Probabilidad"}\NormalTok{)}
\end{Highlighting}
\end{Shaded}

\pandocbounded{\includegraphics[keepaspectratio]{Tallerprobabilidad/figure-latex/unnamed-chunk-11-1.pdf}}

\begin{Shaded}
\begin{Highlighting}[]
\FunctionTok{plot}\NormalTok{(valores\_W, prob\_acumulada, }\AttributeTok{type =} \StringTok{"s"}\NormalTok{,}
     \AttributeTok{main =} \StringTok{"Distribución Acumulada"}\NormalTok{,}
     \AttributeTok{xlab =} \StringTok{"W"}\NormalTok{,}
     \AttributeTok{ylab =} \StringTok{"F(W)"}\NormalTok{)}
\end{Highlighting}
\end{Shaded}

\pandocbounded{\includegraphics[keepaspectratio]{Tallerprobabilidad/figure-latex/unnamed-chunk-12-1.pdf}}

\subsection{Medidas estadísticas}\label{medidas-estaduxedsticas}

\begin{Shaded}
\begin{Highlighting}[]
\NormalTok{valor\_esperado }\OtherTok{\textless{}{-}} \FunctionTok{sum}\NormalTok{(valores\_W }\SpecialCharTok{*}\NormalTok{ probabilidades)}
\NormalTok{varianza }\OtherTok{\textless{}{-}} \FunctionTok{sum}\NormalTok{((valores\_W }\SpecialCharTok{{-}}\NormalTok{ valor\_esperado)}\SpecialCharTok{\^{}}\DecValTok{2} \SpecialCharTok{*}\NormalTok{ probabilidades)}
\NormalTok{desviacion }\OtherTok{\textless{}{-}} \FunctionTok{sqrt}\NormalTok{(varianza)}

\NormalTok{valor\_esperado}
\end{Highlighting}
\end{Shaded}

\begin{verbatim}
## [1] 7
\end{verbatim}

\begin{Shaded}
\begin{Highlighting}[]
\NormalTok{varianza}
\end{Highlighting}
\end{Shaded}

\begin{verbatim}
## [1] 5.833333
\end{verbatim}

\begin{Shaded}
\begin{Highlighting}[]
\NormalTok{desviacion}
\end{Highlighting}
\end{Shaded}

\begin{verbatim}
## [1] 2.415229
\end{verbatim}

\subsection{Simulación}\label{simulaciuxf3n-1}

\begin{Shaded}
\begin{Highlighting}[]
\FunctionTok{set.seed}\NormalTok{(}\DecValTok{123}\NormalTok{)}
\NormalTok{n }\OtherTok{\textless{}{-}} \DecValTok{100000}

\NormalTok{d1 }\OtherTok{\textless{}{-}} \FunctionTok{sample}\NormalTok{(}\DecValTok{1}\SpecialCharTok{:}\DecValTok{6}\NormalTok{, n, }\AttributeTok{replace =} \ConstantTok{TRUE}\NormalTok{)}
\NormalTok{d2 }\OtherTok{\textless{}{-}} \FunctionTok{sample}\NormalTok{(}\DecValTok{1}\SpecialCharTok{:}\DecValTok{6}\NormalTok{, n, }\AttributeTok{replace =} \ConstantTok{TRUE}\NormalTok{)}

\NormalTok{sumas }\OtherTok{\textless{}{-}}\NormalTok{ d1 }\SpecialCharTok{+}\NormalTok{ d2}

\FunctionTok{mean}\NormalTok{(sumas)}
\end{Highlighting}
\end{Shaded}

\begin{verbatim}
## [1] 6.99693
\end{verbatim}

\begin{Shaded}
\begin{Highlighting}[]
\FunctionTok{var}\NormalTok{(sumas)}
\end{Highlighting}
\end{Shaded}

\begin{verbatim}
## [1] 5.806459
\end{verbatim}

\begin{Shaded}
\begin{Highlighting}[]
\FunctionTok{sd}\NormalTok{(sumas)}
\end{Highlighting}
\end{Shaded}

\begin{verbatim}
## [1] 2.409659
\end{verbatim}

\section{Conclusiones}\label{conclusiones}

\begin{itemize}
\tightlist
\item
  Los resultados teóricos coinciden con las simulaciones.
\item
  En combinatoria se aplican coeficientes binomiales.
\item
  En los dados, 7 es el valor más probable.
\end{itemize}

\end{document}
